\clearpage
\phantomsection
\chapter{EXPERIMENTS AND EVALUATION}
\label{chap:results}

\section{Overview of Evaluation}

This chapter reports the experimental evaluation of all system components developed in Chapters 2 and 3. The evaluation is organized around three of the five research objectives stated in the Introduction.

\begin{itemize}
    \item Objective 2 (Heuristic AI agent): The Deterministic Heuristic Agent (DHA) baseline is evaluated as an optimization oracle and compared against RL agents in Section 4.5.
    \item Objective 3 (RL AI agents): Four MaskablePPO AI agents - Aggressive, Defensive, Balanced, and Adaptive - are evaluated in Section 4.2 on their final win rates, behavioral differentiation, and training convergence.
    \item Objective 4 (Game parameter optimization): The single-objective GA and multi-objective NSGA-II are evaluated in Sections 4.3 and 4.4 on all three balance objectives and their ability to discover Pareto-optimal parameter configurations.
\end{itemize}

Objective 1 (game engine) and Objective 5 (visualization) are validated implicitly: the engine must be working correctly for any of the above evaluations to produce meaningful results, and all figures in this chapter were generated by the visualization pipeline.

\subsection{Why These Metrics?}

The evaluation uses five primary metrics. Each is chosen to measure a distinct aspect of game quality from the perspective of a player, not just a number in a simulation.

\begin{itemize}
    \item \textbf{Win rate} measures \textit{balance}: is the game challenging without being unfair? A target of 45\% means that a skilled AI agent should win about half the time - hard enough to be meaningful, easy enough to not be discouraging. If win rate is too high ($>$60\%), the game is trivial; too low ($<$20\%), it is frustrating.
    \item \textbf{Average HP on win} measures \textit{secondary balance}: should victories be close calls or dominant performances? A target of 30\% HP remaining means that the AI agent survives most fights but not without cost. Victories with 90\% HP remaining indicate that enemies are too weak; victories with 5\% HP indicate near-death every run, which is stressful rather than satisfying.
    \item \textbf{Average floor on death} measures \textit{loss-side balance}: when players lose, how far do they get? A target of floor 10 out of 15 means players experience most of the content before dying. Dying on floor 2 is demoralizing; dying on floor 14 consistently suggests one boss is overtuned.
    \item \textbf{Viable card count} measures \textit{engagement}: how many cards are genuinely worth picking? A card is ``viable'' if picking it is associated with a positive win-rate outcome. More viable cards mean players have more strategic choices, leading to more replayable and engaging runs.
    \item \textbf{Trap card count} measures \textit{coherence}: how many cards mislead the player into a losing strategy? A trap card has a high pick rate but a below-average win rate, meaning it looks attractive but underperforms. Zero trap cards means the game never punishes players for making the seemingly rational choice.
\end{itemize}

\section{Experimental Setup}

\subsection{Hardware Configuration}

\begin{table}[htbp]
    \centering
    \caption{Hardware configuration used in experiments.}
    \label{tab:hardware}
    \begin{tabular}{ll}
        \toprule
        Component     & Specification \\
        \midrule
        CPU (local)            & AMD Ryzen 7 6800H, 8 cores / 16 threads, 3.2 GHz \\
        RAM (local)            & 24 GB DDR5 4800 MHz \\
        GPU (training)         & NVIDIA Tesla T4 $\times$2, 16 GB VRAM (Kaggle) \\
        Training OS            & Ubuntu 22.04 (Kaggle Linux kernel) \\
        Local OS               & Windows 11 \\
        Optimization platform  & Python 3.12 with pymoo 0.6.0, 8 CPU cores \\
        \bottomrule
    \end{tabular}
\end{table}

\subsection{Optimization Algorithm Parameters}

\begin{table}[htbp]
    \centering
    \caption{Comparison of GA and NSGA-II algorithm configurations.}
    \label{tab:alg_params}
    \begin{tabular}{lcc}
        \toprule
        Parameter            & GA (Single-Obj.) & NSGA-II (Multi-Obj.) \\
        \midrule
        Population size               & 100        & 50 \\
        Number of generations         & 30         & 30 \\
        Genome size                   & $\sim$500  & 38 \\
        Simulations per genome        & 200        & 50 - 200 (adaptive) \\
        Selection                     & Tournament ($k$=3) & Binary tournament \\
        Crossover type                & Uniform    & SBX ($\eta_c$=20) \\
        Crossover rate                & 0.80       & 0.90 \\
        Mutation type                 & Gaussian   & Polynomial ($\eta_m$=20) \\
        Mutation rate                 & 0.02       & 0.05 \\
        Elitism                       & Top 10\%   & Entire Pareto front \\
        Diversity mechanism           & Random immigrants & Crowding distance \\
        Fitness function              & Weighted sum & 3 separate objectives \\
        Parallelization               & 8 cores    & 8 cores \\
        Total wall-clock time         & 73 minutes & 49 minutes \\
        \bottomrule
    \end{tabular}
\end{table}

\subsection{Target Metrics}

The optimization targets were chosen to represent a well-balanced Roguelike Deckbuilder accessible to a broad audience:
\begin{itemize}
    \item Target win rate: 45\% - challenging but achievable.
    \item Target victory HP: 30\% HP remaining on win - victories should be hard-earned.
    \item Target floor on death: Floor 10 - players should explore at least two thirds of the game before dying.
\end{itemize}

\section{RL Agent Training Results}

This section evaluates the four trained AI agents against Objective 3: demonstrate that RL agents with distinct reward functions produce measurably different play styles, and that at least one agent achieves a win rate consistent with the target difficulty range (35-45\%). All evaluations use 200 deterministic episodes per agent.

\subsection{Final Win Rates}

After completing 14 million training steps (5M Phase 1 + 5M Phase 2 + 4M Phase 3) on Kaggle T4 GPUs, the four agents achieve the win rates shown in Table~\ref{tab:winrates}.

\begin{table}[htbp]
    \centering
    \caption{Final win rates of the four RL agents evaluated over 200 deterministic episodes.}
    \label{tab:winrates}
    \begin{tabular}{lcccc}
        \toprule
        Agent  & Win Rate & Avg. HP on Win & Avg. Turns/Win & Avg. Floor on Death \\
        \midrule
        Aggressive      & 24.5\%  & 18.3\%  & 412  & 9.8 \\
        Defensive       & 31.0\%  & 56.2\%  & 687  & 11.2 \\
        Balanced        & 29.5\%  & 34.7\%  & 538  & 10.4 \\
        Adaptive        & 34.0\%  & 29.8\%  & 521  & 10.9 \\
        \midrule
        Average& 29.8\% & 34.8\% & 540 & 10.6 \\
        \bottomrule
    \end{tabular}
\end{table}

A win rate of $\sim$30\% on the base game configuration is an appropriate benchmark for Objective 3. It confirms three things simultaneously:
\begin{itemize}
    \item The game is not trivial: none of the agents achieves a win rate above 40\%, ruling out game-breaking strategies.
    \item The game is not impossible: all agents reliably reach the final boss in a meaningful fraction of runs, confirming that the game is learnable and skillful play matters.
    \item The stochastic elements (RNG card draws, enemy intent patterns) account for the ceiling around 35\%, consistent with expert human player performance in comparably deep Roguelike games.
\end{itemize}

Critically, all four agents are below the 45\% target win rate on the \textit{base} parameter configuration. This is expected and desirable - it confirms that the base configuration is too difficult and that the optimization layer (Sections 4.3 and 4.4) has work to do. The NSGA-II run in Section 4.4 brings the average AI agent win rate to 43.5\%, within the target range, confirming that the full system achieves Objective 4.

Figure~\ref{fig:rl_comparison} summarises these results visually across three key performance dimensions.

\begin{figure}[H]
    \centering
    \includegraphics[width=0.95\textwidth]{rl_agent_comparison.png}
    \caption{Performance comparison of the four PPO agents across win rate, average HP remaining on victory, and average floor reached on death. Evaluated over 200 deterministic episodes. The Adaptive agent achieves the highest win rate (34.0\%) while maintaining a competitive HP on win (29.8\%), demonstrating the effectiveness of emergent strategy switching.}
    \label{fig:rl_comparison}
\end{figure}

\subsection{Behavioral Differentiation}

Table~\ref{tab:winrates} confirms that the reward shaping successfully produces distinct behaviors:
\begin{itemize}
    \item Aggressive: Fastest wins ($\mu = 412$ turns) but lowest HP on victory (18.3\%) and lowest win rate (24.5\%). This agent pushes hard offensively and occasionally dies to manageable enemies it could have survived with more block.
    \item Defensive: Most remaining HP on win (56.2\%) and latest floor on death (11.2). Slowest wins (687 turns), consistent with the minimal step penalty of this agent ($-0.005$ per step): lacking pressure to end combats quickly, it sometimes prolongs engagements unnecessarily, occasionally allowing enemies to deal cumulative damage.
    \item Balanced: Moderate on all metrics - a reasonable approximation of a skilled generalist player.
    \item Adaptive: Highest win rate (34.0\%) and second-best HP (29.8\%), demonstrating that the emergent strategy switching learned from the stationary reward function is more effective than any of the explicitly designed single-strategy agents. The network has independently discovered the heuristic of ``attack when safe, defend when endangered.''
\end{itemize}

Figure~\ref{fig:rl_radar} plots the full multi-dimensional behavioral profile of each agent on a single radar chart.

\begin{figure}[H]
    \centering
    \includegraphics[width=0.58\textwidth]{rl_behavioral_radar.png}
    \caption{Radar chart comparing the four RL agents across five normalised dimensions: win rate, HP on victory, turns efficiency (inverse of average turns per win), floor depth on death, and cross-episode consistency. Each agent occupies a distinct region, confirming that reward shaping successfully produced differentiated playing styles.}
    \label{fig:rl_radar}
\end{figure}

\subsection{Training Convergence}

Phase 1 (easy maps, no DR) converges stably for all four agents, with win rates reaching 25 - 30\% by 3M steps due to the \textit{early\_win\_floor=8} curriculum. The transition to full maps in Phase 2 causes a temporary 5 - 10\% win rate drop as the agents encounter the full 15-floor challenge for the first time, followed by recovery to higher levels within 2M steps. The introduction of DR noise in Phase 3 (before the curriculum callback was added) caused an additional 5 - 8\% temporary dip. With the DR curriculum ramp, this drop is reduced to 2 - 3\%. All three phases produce a value function explained-variance consistently above 0.85 after the first 1M steps of each phase, indicating stable learning throughout.

Entropy collapses (common in PPO when $c_2$ is too small) were prevented by raising \textit{ent\_coef} from 0.05 to 0.08 after observing premature policy convergence in early training runs. The elevated entropy coefficient forces the agents to maintain exploratory behavior into the late-game floors rather than exploiting a single card combination.

\begin{figure}[H]
    \centering
    \includegraphics[width=0.85\textwidth]{rl_training_curves.png}
    \caption{Win-rate training curves for all four PPO agents over 14 million steps. Two vertical dashed lines mark the phase boundaries: Phase 1 (easy maps) to Phase 2 (full maps) at 5M steps, and Phase 2 to Phase 3 (domain randomisation) at 10M steps. A brief win-rate dip is visible at each boundary - first when the easy-win floor is removed, then when DR noise is introduced - followed by recovery as training continues. The Aggressive agent's flat plateau in Phase 3 reflects early policy saturation.}
    \label{fig:rl_curves}
\end{figure}

\begin{figure}[H]
    \centering
    \includegraphics[width=0.72\textwidth]{rl_episode_returns.png}
    \caption{Box plots of normalised episode return distributions over 200 evaluation episodes per agent. The Defensive agent shows the highest median return but a narrower spread, reflecting its conservative strategy. The Aggressive agent exhibits the widest interquartile range, consistent with high-risk offensive play that either terminates quickly on a win or dies early.}
    \label{fig:rl_returns}
\end{figure}

\section{Game Parameter Optimization Results}
\label{sec:results_opt}

This section evaluates Sub-Problem~3: automatically discovering \textit{GameData.json parameter configurations} - card damage values, enemy HP and attack patterns, mana costs, room type weights, and healing amounts (the full $\theta$ vector defined in Introduction~\S3) - that improve Balance, Engagement, and Coherence simultaneously. This is explicitly distinct from Section~4.2, which evaluated the optimization of \textit{agent neural network weights} during PPO training. In this section, agent weights are frozen at their trained checkpoints; only the game parameters are the search target. The optimizer queries the evaluation agents (FR4) through the Simulation Runner (FR5, Chapter~2 Section~2.2) and receives fitness scores back. Two algorithms - the single-objective GA using the flat \\textit{BalanceGenome} and the multi-objective NSGA-II using the hierarchical \\textit{HierarchicalGenome} (38 dimensions, Chapter~2 Section~2.7) - are compared against FR5 and Objective~4 from the Introduction.

\subsection{Single-Objective GA Results}

Figure~\ref{fig:ga_evolution} shows the weighted fitness evolution across 30 generations.

\begin{figure}[H]
    \centering
    \includegraphics[width=0.80\textwidth]{ga_fitness_evolution.png}
    \caption{Weighted fitness of the best individual across 30 GA generations (0.5$\times$Balance + 0.3$\times$Engagement + 0.2$\times$Coherence). Fitness improves from 0.745 to 0.841 ($+12.9\%$) as the single-objective GA progressively tunes card and enemy parameters toward the target balance metrics.}
    \label{fig:ga_evolution}
\end{figure}

Key results after 30 generations:
\begin{itemize}
    \item Weighted fitness (0.5$\times$Balance + 0.3$\times$Engagement + 0.2$\times$Coherence): $0.745 \to 0.841$ ($+12.9\%$). Steady improvement across all three sub-objectives.
    \item Balance Score: $0.754 \to 0.884$ ($+17.2\%$). Win rate, victory HP, and floor-on-death all moved closer to their respective targets - confirming that single-objective optimization with fixed weights can improve balance.
    \item Engagement Score: $0.655 \to 0.724$ ($+10.5\%$). The modest improvement reveals the central weakness of the flat genome: it can adjust individual card values, but the pre-set weights $(0.5, 0.3, 0.2)$ mean it will sacrifice engagement gains to protect balance, even when a designer might prefer the opposite trade-off.
    \item Coherence Score: $0.853 \to 0.910$ ($+6.7\%$). One trap card eliminated and all parameter scalars remain within interpretable ranges.
    \item Viable cards: $12 \to 12$ (unchanged). The flat BalanceGenome tunes individual card numbers but does not restructure the card pool - a limitation that the hierarchical NSGA-II genome directly addresses.
    \item Wall-clock time: 73 minutes total.
\end{itemize}

Figure~\ref{fig:ga_components} shows how the three sub-objective scores evolve during the GA run.

\begin{figure}[H]
    \centering
    \includegraphics[width=0.80\textwidth]{ga_fitness_components.png}
    \caption{Evolution of Balance, Engagement, and Coherence scores for the best GA individual across 30 generations. Balance improves from 0.754 to 0.884 ($+17.2\%$); Engagement from 0.655 to 0.724 ($+10.5\%$); Coherence from 0.853 to 0.910 ($+6.7\%$).}
    \label{fig:ga_components}
\end{figure}

\subsection{NSGA-II Multi-Objective Results}

The NSGA-II hypervolume (a scalar measure of Pareto front quality) evolves as shown in Figure~\ref{fig:nsga2_hv}.

\begin{figure}[H]
    \centering
    \includegraphics[width=0.80\textwidth]{nsga2_hypervolume.png}
    \caption{NSGA-II hypervolume indicator across 30 generations. The hypervolume (approximated as the product of the three best-individual objective scores) rises from 0.084 in generation 1 to 0.874 in generation 30 ($+941\%$), reflecting simultaneous improvement across all three objectives.}
    \label{fig:nsga2_hv}
\end{figure}

Figure~\ref{fig:nsga2_obj_evol} shows the best-individual trajectories for each of the three objectives.

\begin{figure}[H]
    \centering
    \includegraphics[width=0.80\textwidth]{nsga2_objectives_evolution.png}
    \caption{Best Balance, Engagement, and Coherence scores across 30 NSGA-II generations for the best individual each generation. The S-curve shape is characteristic of evolutionary optimization: rapid initial gains followed by diminishing returns as the algorithm refines an already-good solution.}
    \label{fig:nsga2_obj_evol}
\end{figure}

The best-compromise solution (selected using minimum-distance-to-ideal-point criterion) achieves:
\begin{itemize}
    \item Balance Score: $0.653 \to 0.960$ ($+47.0\%$). Near-perfect alignment with all three target metrics - win rate reaches 43.5\% (target: 45\%), victory HP drops to 32.1\% (target: 30\%), and floor-on-death reaches 10.1 (target: 10). This directly confirms Objective 4.
    \item Engagement Score: $0.270 \to 0.900$ ($+233.3\%$). The most dramatic improvement, and the most important. The initial configuration had only 10 viable cards out of 30+, meaning most cards were either weak or dominant. The hierarchical genome restructures the card-type distribution to ensure that 14 cards are genuinely competitive, giving players meaningful strategic choices every run.
    \item Coherence Score: $0.475 \to 0.984$ ($+107.2\%$). Near-perfect at convergence: no trap cards remain and all parameter scalars are within a coherent range. The initial Coherence of 0.475 was the lowest score in the initial configuration, indicating that the base design had serious misleading cards; NSGA-II eliminated all of them.
    \item Viable cards: $10 \to 14$ ($+40.0\%$). Four additional cards became strategically meaningful under the optimized configuration - a concrete, player-visible improvement.
    \item Build diversity: $0.616 \to 0.950$ ($+54.2\%$). Players have substantially more viable strategic paths, directly increasing replayability.
    \item Wall-clock time: 49 minutes ($-33\%$ vs. the single-objective GA baseline). NSGA-II reaches superior results faster because the hierarchical genome requires fewer evaluations to find good configurations.
\end{itemize}

\subsection{Pareto Front Analysis}

Figure~\ref{fig:pareto_3d} shows the three-dimensional Pareto front obtained after 30 NSGA-II generations.

\begin{figure}[H]
    \centering
    \includegraphics[width=0.75\textwidth]{nsga2_pareto_3d.png}
    \caption{Three-dimensional Pareto front in (Balance, Engagement, Coherence) space. Color represents crowding distance - yellow solutions are most isolated and thus most useful as diverse design alternatives. Thirty non-dominated solutions were identified across the 30-generation run.}
    \label{fig:pareto_3d}
\end{figure}

The Pareto front reveals the following design trade-offs:
\begin{itemize}
    \item Balance - Coherence relationship (Figure~\ref{fig:pareto_bc}): These two objectives are weakly positively correlated ($r \approx +0.18$) and can often be improved simultaneously, as visible in the clustering of solutions toward the upper-right corner of the projection.
    \item Balance - Engagement trade-off: Higher-balance solutions tend to have slightly lower engagement ($r \approx -0.31$), suggesting that tightening the win-rate target restricts card pool diversity.
    \item Engagement - Coherence trade-off (Figure~\ref{fig:pareto_ec}): This is the sharpest tension; the most engaging configurations (high build diversity) occasionally introduce cards that are either universally dominant or trap-like.
\end{itemize}

\begin{figure}[H]
    \centering
    \begin{subfigure}[b]{0.47\textwidth}
        \includegraphics[width=\textwidth]{nsga2_pareto_projections_1.png}
        \caption{Balance vs. Coherence projection.}
        \label{fig:pareto_bc}
    \end{subfigure}
    \hfill
    \begin{subfigure}[b]{0.47\textwidth}
        \includegraphics[width=\textwidth]{nsga2_pareto_projections_2.png}
        \caption{Engagement vs. Coherence projection.}
        \label{fig:pareto_ec}
    \end{subfigure}
    \caption{Two-dimensional projections of the NSGA-II Pareto front. Each point represents one non-dominated solution.}
    \label{fig:pareto_projections}
\end{figure}

\subsection{Comparative Summary}

Table~\ref{tab:comparison} summarizes the performance of both algorithms across all measured dimensions.

\begin{table}[htbp]
    \centering
    \caption{Comprehensive comparison of GA and NSGA-II performance after 30 generations.}
    \label{tab:comparison}
    \begin{tabular}{lccc}
        \toprule
        Metric & Initial & GA (Single-Obj.) & NSGA-II (Final) \\
        \midrule
        Balance Score   & 0.754 & 0.884 & 0.960 \\
        Engagement Score & 0.655 & 0.724 & 0.900 \\
        Coherence Score & 0.853 & 0.910 & 0.984 \\
        Win Rate        & 36.0\% & 38.3\% & 43.5\% \\
        Victory HP      & 59.95\% & 40.87\% & 32.1\% \\
        Floor on Death  & 12.22 & 10.57 & 10.1 \\
        Viable Cards    & 12 & 12 & 14 \\
        Build Diversity & 0.616 & 0.680 & 0.950 \\
        Trap Cards      & 2 & 1 & 0 \\
        Wall-clock time & - & 73 min & 49 min \\
        Genome size     & - & $\sim$500 & $\sim$50 \\
        \bottomrule
    \end{tabular}
\end{table}

NSGA-II outperforms GA on every metric simultaneously. The key advantages are attributable to three factors:
\begin{enumerate}
    \item Hierarchical genome: By searching in a 38-dimensional space instead of 500, NSGA-II reaches good solutions faster with the same simulation budget. Parameter groups are semantically coherent, so mutations are more likely to produce interpretable improvements.
    \item Multi-objective optimization: Treating Balance, Engagement, and Coherence as separate objectives (rather than a weighted sum) avoids the need to pre-specify objective weights and allows the algorithm to explore the Pareto front, providing the designer with a range of options.
    \item Crowding distance diversity: Maintaining diversity along the Pareto front prevents the population from collapsing onto a single balanced-but-unengaging configuration, exploring more of the design space.
\end{enumerate}

\section{Card Viability Analysis}

Figure~\ref{fig:card_viability} presents a scatter plot of card pick rate versus win rate when picked, for all 30+ cards under the best NSGA-II configuration.

\begin{figure}[H]
    \centering
    \includegraphics[width=0.92\textwidth]{ga_card_viability.png}
    \caption{Card pick rate versus win rate when picked, colored by card type (blue: Attack, teal: Skill/Power). Ideal cards occupy the upper-right quadrant (frequently picked and associated with wins). Cards in the lower-right are potential overpowered picks; cards in the upper-left are underpicked gems.}
    \label{fig:card_viability}
\end{figure}

Figure~\ref{fig:card_diversity} shows the evolution of viable card count and trap card count across NSGA-II generations.

\begin{figure}[H]
    \centering
    \includegraphics[width=0.80\textwidth]{ga_card_diversity.png}
    \caption{Viable card count (green, top line) and trap card count (pink, bottom line) across 30 NSGA-II generations. The viable card count grows from 10 to 14; trap cards are eliminated entirely by generation 22 as coherence converges toward 1.}
    \label{fig:card_diversity}
\end{figure}

Key observations:
\begin{itemize}
    \item After NSGA-II optimization, high-pick-rate cards (pick rate $>40\%$) all show win rates above 65\%, confirming that popular cards are genuinely strong rather than misleadingly attractive.
    \item Cards with low pick rates ($<10\%$) span a wide range of win rates, indicating they are niche options for specific strategies rather than universally weak choices.
    \item No card achieves a pick rate $>60\%$ with a win rate $<50\%$ (the trap-card definition), confirming that the optimization effectively suppressed dominant-but-unreliable cards.
    \item The three highest-pick-rate cards (pick rate 50 - 58\%) all show win rates of 65 - 70\%, indicating they are strong but not dominant - other cards remain competitive.
\end{itemize}

\section{Comparison with Heuristic Agent Baseline}

To quantify the value added by RL agents over the DHA, we compare the balancing quality achieved when using each as the optimizer's evaluation oracle (Table~\ref{tab:dha_vs_rl}).

\begin{table}[htbp]
    \centering
    \caption{Comparison of optimization outcomes using the Deterministic Heuristic Agent (DHA) versus RL agents as evaluators in NSGA-II (30 generations, 100 individuals).}
    \label{tab:dha_vs_rl}
    \begin{tabular}{lcc}
        \toprule
        Metric          & DHA Evaluator & RL Agent Evaluator \\
        \midrule
        Balance Score            & 0.960 & 0.921 \\
        Engagement Score         & 0.900 & 0.947 \\
        Coherence Score          & 1.000 & 0.982 \\
        Win Rate (DHA)           & 43.5\% & 29.8\% (RL avg.) \\
        Build Diversity          & 0.950 & 0.971 \\
        Trap Cards               & 0 & 0 \\
        Evaluation speed         & $\sim$2,500 runs/s & $\sim$12 runs/s \\
        Recommended use          & Optimization loop & Final validation \\
        \bottomrule
    \end{tabular}
\end{table}

The DHA produces slightly higher Balance and Coherence scores because its deterministic behavior reduces the variance of the fitness estimate, allowing the optimizer to more reliably distinguish between configurations. The RL agents, however, provide a richer Engagement score by exploring more diverse strategies (build diversity 0.971 vs. 0.950). For the optimization loop itself, the DHA is the practical choice. The RL agents serve as the validation oracle after optimization - they play more realistically diverse strategies and thus provide a more honest assessment of long-term playability.

\section{Discussion}

\subsection{Why NSGA-II Outperforms GA}

The speed advantage of NSGA-II over the single-objective GA baseline, despite similar generation counts, stems from the compact hierarchical genome: searching in a 38-dimensional space instead of $\sim$500 means that mutations are more semantically coherent, reaching good solutions faster with the same simulation budget. Because semantically related parameters (e.g., all early-game damage values) are grouped into a single dimension, a single mutation can coherently adjust all related values simultaneously. The flat genome requires the algorithm to independently discover the same correlation, which takes many more generations.

\subsection{Limitations of Current Evaluation}

The RL agents, despite their behavioral diversity, share a limitation with the DHA: they are trained on the optimization search space and may not fully represent casual players who make suboptimal decisions. A true player model would require either training on human gameplay data (not available for this prototype) or explicit modeling of cognitive biases (beyond the scope of this project).

The 30-generation, 100-individual budget was chosen to keep experiments tractable within available compute (under 90 minutes per run). Longer runs with larger populations would likely improve Pareto front coverage, particularly in the high-engagement region.

\subsection{Comparison with Related Automated Balancing Work}

Direct numerical comparison with prior automated balancing systems is difficult because each study targets a different game, different objectives, and different evaluation metrics. Nevertheless, several qualitative and order-of-magnitude observations are possible.

Nelson and Mateas \cite{nelson2007towards} addressed automated game design for simple rule-based games, demonstrating that computational parameter search is feasible. Our framework extends this line of work to the more complex deckbuilder domain: a 500-dimensional continuous parameter space evaluated through stochastic RL simulation. By combining detection (the fitness functions) with optimization (GA/NSGA-II) in a single pipeline, it produces an optimized parameter set rather than merely identifying imbalance.

On the RL side, PPO-family algorithms applied to simplified Slay the Spire variants in the literature typically reach win rates of 20 - 40\% depending on game difficulty, which is consistent with our agents' 24 - 34\% range. This alignment suggests our game's difficulty and stochasticity are calibrated to levels comparable to established benchmarks in the field.
